\documentclass[../main.tex]{subfiles}
\begin{document}

\chapter{メモリ一貫性}
\lessoninfo{
  video={Lesson 21，動画 1--14},
  textbook={H\&P \cite{hennessy2017} §5.6；\cite{mosberger1993}},
  keywords={メモリ一貫性，キャッシュコヒーレンス，プログラム順序，一貫性モデル，逐次一貫性，強い一貫性，ロードの再実行，緩和一貫性，MSYNC，データ競合，弱い一貫性，リリース一貫性},
}

第19章は，メモリ位置の読み出しがその位置へ最後に書かれた値を返すように
し，第20章はロックとバリアを用いて，共有データを更新するスレッドを協調
させた．どちらも，\emph{異なる}位置へのアクセスがどの順序で効力を生じるか
には触れておらず，アウトオブオーダプロセッサは当然のようにその順序を変え
る．この章では，そのような順序変更が，プログラム順どおりのどの実行でも
生じえない結果を生みうることを示し，そうした結果を排除する保証である
逐次一貫性を定義し，アウトオブオーダプロセッサがそれをどう提供するかを
説明する．続いて，順序が重要な箇所をプログラム自身が明示する緩和一貫性と，
プログラムがその違いを見分けられるかどうかを決めるデータ競合を扱う．

\section{コヒーレンスと一貫性}
\label{sec:l21-coherence}

キャッシュコヒーレンス\index{きゃっしゅこひーれんす@キャッシュコヒーレンス}%
（\cref{sec:l19-definition}）は，1つのメモリ位置へのアクセスを制約する．
すなわち，どのコアによるその位置の読み出しも，どのコアによるものであれ，
その同じ位置への最も新しい書き込みの値を返す．\source{Lesson 21，動画
1--2；H\&P §5.6．}これによってはじめてスレッドはデータを共有できる．しかし
コヒーレンスは各位置を個別に扱う．位置 A に書いてから位置 B を読むコアが，
他のコアから見てこの順序で2つのアクセスを行うかどうかについても，1つの
コアによる A と B への2つの書き込みが，そのコアが行った順序で他のコアに
見えるようになるかどうかについても，何も述べていない．

\begin{definition}[メモリ一貫性]
  共有メモリシステムの\term[めもりいっかんせい]{メモリ一貫性}[memory
  consistency]とは，異なる位置へのアクセスも含め，すべてのコアのメモリ
  アクセスがどの順序で効力を生じるかについての，そのシステムの保証で
  ある．これは，並列プログラムの読み出しが返しうる値の組合せを決める．
\end{definition}

書き込みは，読み出しから見えるようになった時点で効力を生じる．読み出しは，
それが返すメモリ内容が存在した時点で効力を生じる．システムが許す順序を
厳密に述べたものが\term[いっかんせいもでる]{一貫性モデル}[consistency
model]である．この章では最も重要なモデルである逐次一貫性とその緩和版を
展開し，\cref{sec:l21-models} でさらに他のモデルを概観する．

2つの性質の由来は異なる．コヒーレンスが必要なのは，キャッシュがメモリの
コピーを保持するからである．キャッシュのないシステムはどの位置にもコピー
を1つしか持たず，構成上コヒーレントである．一貫性はキャッシュに依存しない．
どの位置にもコピーが1つしかなくても，メモリアクセスをプログラムが指定する
順序と異なる順序で行うプロセッサは，自分の読み出しが返す値を変えてしまう
---そして次節が示すように，現在のプロセッサはまさにそれを行う．したがって，
コヒーレントなシステムが一貫性を持つとは限らない．

\section{一貫性が問題になる理由}
\label{sec:l21-matters}

異なる位置へのアクセスの順序が問題になるかどうかは，小さなプログラム
（\cref{ex:l21-flags}）で最もよく分かる．\source{Lesson 21，動画 3--5；
H\&P §5.6．}

\begin{example}[2つのフラグ]
  \label{ex:l21-flags}
  2つのコアが同じコードを実行する．コア 0 では R4 がフラグ X のアドレス
  を，R5 がフラグ Y のアドレスを保持し，コア 1 では X と Y の役割が
  入れ替わる．どちらのフラグも初めは 0 である．
  \begin{lstlisting}[language={[x86masm]Assembler}, numbers=none]
      ADDI R1, R0, 1    ; R1 <- 1
      SW   R1, 0(R4)    ; 自分のフラグ <- 1
      LW   R2, 0(R5)    ; R2 <- 相手のフラグ
      BNEZ R2, skip     ; 相手のコアが宣言済み
      ...               ; 片方のコアだけが行う作業
skip: ...
\end{lstlisting}
  各コアは，これから作業を行うことをまず自分のフラグで宣言し，相手のコア
  が宣言していない場合にだけ作業を行う．
\end{example}

各コアが2つのメモリアクセスをプログラム順に行えば，作業を行うコアは多くとも
1つである．コア 0 が Y を 0 と読んだとする．するとコア 0 の Y の読み出しは
コア 1 の Y へのストアより前に効力を生じており，したがって，そのストアに
続くコア 1 の X の読み出しよりも前である．コア 0 の X へのストアは自分の
Y の読み出しに先行するので，コア 1 の X の読み出しにも先行し，その読み出し
は 1 を返す．（両方のストアが両方の読み出しより前に効力を生じれば，どちらの
コアも作業を行わない．）\margin{2つのフラグのプログラムは，
通常のロードとストアを使う Dekker や Peterson の相互排除アルゴリズム
（第20章）の核心である．}

アウトオブオーダプロセッサはアクセスをプログラム順に行わない．ストアが
メモリに書くのはコミットするときであるのに対し，ロードは，アドレスを計算
し，ロードストアキュー（\cref{sec:l09-lsq}）に同じアドレスへの先行する
ストアがないと分かり次第メモリを読む．\cref{ex:l21-flags} のストアとロード
は異なるアドレスにアクセスするので，ロードはストアがコミットするよりはるか
前にメモリを読みうる．\margin{その距離はリオーダバッファで決まる．Skylake コ
アでは 224 エントリである（第8章）．Agner Fog，
マイクロアーキテクチャマニュアル．}これが両方のコアで起これば
（\cref{fig:l21-reorder}），
両方のロードが 0 を読み，両方のコアが作業を行う．

\begin{figure}[htbp]
  \centering
  % Two out-of-order cores run the two-flag program.  In program order each
% core stores 1 into its own flag and then loads the other core's flag.
% Both loads read memory before either store commits, so both read 0.
% Arrows: data read from memory (up/down towards a core) and stores written
% to memory at commit (towards the memory row).
\begin{tikzpicture}[x=0.95cm,y=1cm, font=\footnotesize\sffamily, >={Stealth[length=1.6mm]},
  row/.style={ln-muted, line width=0.6pt},
  ev/.style={circle, fill=black, inner sep=1.3pt},
  lab/.style={font=\scriptsize, align=center, inner sep=1pt},
  mem/.style={draw, fill=ln-box, minimum height=0.5cm, inner sep=1pt, font=\scriptsize}]
  % ---- row names
  \node[anchor=east] at (-0.15,1.5) {\iflangja{コア}{core}~0};
  \node[anchor=east] at (-0.15,0) {\iflangja{メモリ}{memory}};
  \node[anchor=east] at (-0.15,-1.5) {\iflangja{コア}{core}~1};
  \draw[row] (0,1.5) -- (10.2,1.5);
  \draw[row] (0,-1.5) -- (10.2,-1.5);
  % ---- memory contents over time
  \node[mem, minimum width=5.0*0.95cm, anchor=west] at (0,0) {X $=$ 0, Y $=$ 0};
  \node[mem, minimum width=2.2*0.95cm, anchor=west] at (5.0,0) {X $=$ 1, Y $=$ 0};
  \node[mem, minimum width=3.0*0.95cm, anchor=west] at (7.2,0) {X $=$ 1, Y $=$ 1};
  % ---- core 0: labels above its row
  \node[ev] at (0.7,1.5) {};
  \node[lab, anchor=south] at (0.7,1.62) {\iflangja{SW X\\実行}{SW X\\executes}};
  \node[ev, ln-caution] at (2.4,1.5) {};
  \node[lab, anchor=south, text=ln-caution] at (2.4,1.62) {\iflangja{LW Y\\0 を読む}{LW Y\\reads 0}};
  \draw[->, ln-caution] (2.4,0.27) -- (2.4,1.4);
  \node[ev] at (5.0,1.5) {};
  \node[lab, anchor=south] at (5.0,1.62) {\iflangja{SW X\\コミット}{SW X\\commits}};
  \draw[->] (5.0,1.4) -- (5.0,0.27);
  \node[ev, ln-caution] at (8.4,1.5) {};
  \node[lab, anchor=south, text=ln-caution] at (8.4,1.62) {\iflangja{LW Y コミット，\\作業を実行}{LW Y commits,\\does the work}};
  % ---- core 1: labels below its row
  \node[ev] at (1.3,-1.5) {};
  \node[lab, anchor=north] at (1.3,-1.62) {\iflangja{SW Y\\実行}{SW Y\\executes}};
  \node[ev, ln-caution] at (3.6,-1.5) {};
  \node[lab, anchor=north, text=ln-caution] at (3.6,-1.62) {\iflangja{LW X\\0 を読む}{LW X\\reads 0}};
  \draw[->, ln-caution] (3.6,-0.27) -- (3.6,-1.4);
  \node[ev] at (7.2,-1.5) {};
  \node[lab, anchor=north] at (7.2,-1.62) {\iflangja{SW Y\\コミット}{SW Y\\commits}};
  \draw[->] (7.2,-1.4) -- (7.2,-0.27);
  \node[ev, ln-caution] at (9.4,-1.5) {};
  \node[lab, anchor=north, text=ln-caution] at (9.4,-1.62) {\iflangja{LW X コミット，\\作業を実行}{LW X commits,\\does the work}};
  % ---- time axis
  \draw[->, ln-muted] (0,-2.45) -- (10.2,-2.45)
    node[below left, inner sep=2pt, font=\scriptsize] {\iflangja{時間}{time}};
\end{tikzpicture}

  \caption{2つのアウトオブオーダコア上での \cref{ex:l21-flags}．各コアの
    ロードは自分のストアがコミットする前にメモリを読み，両方のロードが 0
    を読んで，両方のコアが作業を行う．}
  \label{fig:l21-reorder}
\end{figure}

この実行ではどの読み出しもコヒーレントである．各ロードが相手のフラグを
読む時点で，そのフラグは実際に 0 を保持しており，ストアはコミット時に，
プログラム順に効力を生じる．\caution{ここではどのキャッシュも古い値を保持し
ていないので，どのコヒーレンスプロトコルもこの結果を防げない．
原因は古いデータではなく順序変更である．}それでもこの結果は，
プログラム順どおりの
どの実行でも生じえない．アウトオブオーダのコアが変えたのは，異なる位置への
アクセスの順序である．すなわち，各コアによる相手のフラグの読み出しが，
自分のストアより前に効力を生じたのである．

\subsection{正しいプログラムにおけるロードの順序変更}
\label{sec:l21-ready}

2つのフラグのプログラムは作為的に見えるかもしれないが，同じ種類の順序変更
は，よくある，それ自体は正しいパターンを壊す．その1つが\emph{データ準備
完了フラグ}である．書き込み側スレッドは新しいデータをストアしてからフラグ
を立て，読み出し側スレッドはフラグが立つまで待ってからデータを使う．R4 が
フラグのアドレスを，R5 がデータのアドレスを保持するとき，書き込み側は
\begin{lstlisting}[language={[x86masm]Assembler}, numbers=none]
      SW   R6, 0(R5)    ; data  <- 新しい値
      SW   R7, 0(R4)    ; ready <- R7 (= 1)
\end{lstlisting}
を，読み出し側は
\begin{lstlisting}[language={[x86masm]Assembler}, numbers=none]
poll: LW   R1, 0(R4)    ; R1 <- ready
      BEQZ R1, poll     ; 未設定なら再度ポーリング
      LW   R2, 0(R5)    ; R2 <- data
\end{lstlisting}
を実行する．プログラム順では，読み出し側のデータのロードは，フラグが立って
いるのを見たフラグのロードに続き，そのロードは書き込み側のフラグへのストア
に続き，そのストアはさらに書き込み側のデータのストアに続く．したがって
読み出し側は常に新しいデータを得る．分岐予測（第4章）はこの保証を失わせる．
読み出し側のコアはループを抜けると予測し，データのロードを投機的に，場合に
よっては同じ反復のフラグのロードより前に実行する．\cref{tab:l21-ready} は，
データが 3 を保持し，書き込み側が 8 をストアする，そのような実行を追った
ものである．フラグのロードは書き込み側がフラグを立てた後に実行されるので，
予測は結果として正しく，データのロードは古い値のままコミットする．

\begin{table}[htbp]
  \centering\small
  \caption{アウトオブオーダの読み出し側コア（コア 0）と書き込み側コア
    （コア 1）におけるデータ準備完了フラグ．各手順の後のメモリ内容．}
  \label{tab:l21-ready}
  \begin{tabular}{@{}cclcc@{}}
    \toprule
    手順 & コア & 事象 & ready & data \\
    \midrule
    1 & 0 & \texttt{LW data} が予測経路で実行される: R2 $=$ 3 & 0 & 3 \\
    2 & 1 & \texttt{SW data} がコミットする                   & 0 & 8 \\
    3 & 1 & \texttt{SW ready} がコミットする                  & 1 & 8 \\
    4 & 0 & \texttt{LW ready} が実行される: R1 $=$ 1          & 1 & 8 \\
    5 & 0 & \texttt{BEQZ} が解決する: 予測どおり不成立        & 1 & 8 \\
    6 & 0 & 3つの命令が R2 $=$ 3 のままコミットする           & 1 & 8 \\
    \bottomrule
  \end{tabular}
\end{table}

あるスレッドが別のスレッドを生成し，そのスレッドが終了するまで待ってから，
それが生成した結果を使う場合にも同じことが起こる．待機が終わるのは，
スレッドの終了時にオペレーティングシステムが書き込む状態をロードし，それが
終了を示した時点である．プログラム順でそれに続く結果のロードは，それより
前に，すなわち他方のスレッドが結果を書く前に実行されうる．
\margin{スレッドライブラリはこれを隠す．POSIX は \texttt{pthread\_join} に，
終了するスレッドの書き込みを join する側のアクセスより前に順序付けることを要
求している．}

どちらのパターンでも，すべてのキャッシュはコヒーレントであり，アクセスが
プログラム順に効力を生じる限りプログラムは正しい．そもそもこのような
プログラムを書けるようにするには，プログラマは異なる位置へのアクセスの
順序についての保証---コヒーレンスを超えた保証---を必要とする．

\begin{keypoint}
  コヒーレンスは各位置へのアクセスを別々に順序付ける．アウトオブオーダ実行
  と分岐予測は，他の位置への先行するアクセスより前にロードが効力を生じる
  ことを許し，その結果，プログラム順どおりのどの実行でも生じえない結果が
  現れうる．
\end{keypoint}

\section{逐次一貫性}
\label{sec:l21-sc}

最も直観的な保証は，コアが順番に交代したかのようにメモリが振る舞うことを
要求する．\source{Lesson 21，動画 6；H\&P §5.6；\cite{lamport1979}．}

\begin{definition}[逐次一貫性]
  共有メモリシステムは，どの実行の結果も，各コアのメモリアクセスがプログラム
  順に行われ，異なるコアのアクセスが何らかの任意の順序でインターリーブ
  された場合と同じであるとき，\term[ちくじいっかんせい]{逐次一貫性}%
  [sequential consistency]を提供するという．
\end{definition}

Lamport \cite{lamport1979} に遡るこの定義には2つの部分がある．1つのコアの
内部では順序は固定されている．すなわち，アクセスはプログラムが指定する順序
で現れる．コアをまたぐインターリーブは任意であり，ある実行がどれに対応する
かはコアのタイミングによるので，正しいプログラムはあらゆるインターリーブ
に対して正しい結果を生まなければならない．「かのように」という言葉が重要
である．どのプログラムも違いを見分けられない限り，ハードウェアが実際に
そのような順序でアクセスを行う必要はない．

\begin{example}[2つのフラグのプログラムのインターリーブ]
  \label{ex:l21-interleavings}
  \cref{ex:l21-flags} では，コア 0 が X をストアし（$w_0$），次に Y を読む
  （$r_0$）．コア 1 は Y をストアし（$w_1$），次に X を読む（$r_1$）．各コア
  の順序を保つと，6 通りのインターリーブが残る
  （\cref{tab:l21-interleavings}）．\tip{両コアのプログラム順を保つインター
  リーブの数は，アクセスが $m$ 個と $n$ 個なら $\binom{m+n}{m}$，
  ここでは $\binom{4}{2}=6$ である．}
  そのいずれにおいても少なくとも一方の読み出しが 1 を返すので，逐次一貫性
  を持つシステムが両方のコアに作業を行わせることはない．すなわち
  \cref{fig:l21-reorder} の実行は逐次一貫性を持たない．
\end{example}

\begin{table}[htbp]
  \centering\small
  \caption{各コアのプログラム順を保つ，\cref{ex:l21-interleavings} の
    アクセスの 6 通りのインターリーブ．}
  \label{tab:l21-interleavings}
  \begin{tabular}{@{}clccl@{}}
    \toprule
    & 順序 & $r_0$（Y） & $r_1$（X） & 作業を行うコア \\
    \midrule
    1 & $w_0\;r_0\;w_1\;r_1$ & 0 & 1 & コア 0 \\
    2 & $w_0\;w_1\;r_0\;r_1$ & 1 & 1 & なし \\
    3 & $w_0\;w_1\;r_1\;r_0$ & 1 & 1 & なし \\
    4 & $w_1\;w_0\;r_0\;r_1$ & 1 & 1 & なし \\
    5 & $w_1\;w_0\;r_1\;r_0$ & 1 & 1 & なし \\
    6 & $w_1\;r_1\;w_0\;r_0$ & 1 & 0 & コア 1 \\
    \bottomrule
  \end{tabular}
\end{table}

\needspace{8\baselineskip}
\begin{note}[強い一貫性]
  逐次一貫性が要求するのは，\emph{何らかの}インターリーブが結果を説明する
  ことだけである．そのインターリーブはアクセスが実際に起こる実時間と一致
  する必要がないので，他のコアがその位置への書き込みをすでに完了していて
  も，読み出しが古い値を返すことがある．順序が実時間にも従わなければ
  ならない---他のアクセスの開始前に完了したアクセスが先になり，どの
  読み出しも最も新しく書かれた値を返す---システムは
  \term[つよいいっかんせい]{強い一貫性}[strong consistency]を提供する．
  Mosberger \cite{mosberger1993} はこれをアトミック一貫性と呼んでいる．
\end{note}

\section{逐次一貫性の実現}
\label{sec:l21-implementing}

\subsection{単純な実現}

逐次一貫性を提供する最も単純な方法は，その定義から「かのように」を取り除く
ことである．\source{Lesson 21，動画 7--8；H\&P §5.6．}各コアはメモリアクセス
を一度に1つずつ，プログラム順に行う．すなわち，先行するアクセスがすべて
完了してはじめて次のアクセスを開始する．コヒーレンスによって，各アクセスは
効力を生じると同時にすべてのコアから見えるようになるので，すべてのコアの
アクセスが効力を生じる順序そのものが，定義の要求する種類のインターリーブ
になる．

この実現は正しいが非常に遅い．アウトオブオーダコアはメモリにアクセスしない
命令の順序を依然として変えられるが，進行中のメモリアクセスは一度に1つしか
持てない．ロードは先行するストアがすべてコミットするまで実行できず，
ロードした値を使う命令はすべてそれとともに待つ．しかしメモリアクセスこそ，
アウトオブオーダ実行が最も重ね合わせる必要のある遅い操作である．キャッシュ
ミスは多くのサイクルを要し，マルチプロセッサではさらに，アクセスがバスと
コヒーレンスプロトコル（第19章）を待たなければならないこともある．
\margin{直列化すると各レイテンシを丸ごと支払う．Skylake 級のコアで L1 ヒッ
ト 4 サイクル，L2 は 12，L3 は 44，メモリは 60--100~ns である（第14章）．}メ
モリ
アクセスが一度に1つに制限されると，コアはアウトオブオーダ実行の利点の大半
を失う．

\subsection{ロードを安全に順序変更する}

より良い実現は，コアにロードをアウトオブオーダで実行させ，それが逐次一貫性
を破りうる場合を検出して修復する．これを可能にする観察が2つある．
\begin{itemize}
  \item \textbf{注意が必要なのはロードだけである．}ストアはコミット時に
    プログラム順にメモリに書く（\cref{sec:l09-write}）ので，すでにプログラム
    順に効力を生じている．同じコアの先行するストアの位置を読むロードは，
    そのストアの値をロードストアキューから得る（第9章）．
  \item \textbf{ROB がプログラム順を記録している．}実行の順序がどうであれ，
    ROB は命令をプログラム順に保持しており，早く実行されたロードは，
    コミットする時点で効力を生じたとみなせる．読んだ値は，その間に他のコア
    がその位置に書いていない限り，その時点においても正しい．
\end{itemize}
そこでコアは，実行済みでまだコミットしていないすべてのロードの位置について，
他のコアのコヒーレンストラフィックを監視する．書き込み時無効化方式のもとで
は，他のコアによる書き込みは，その位置を含むブロックの無効化として届く
（\cref{sec:l19-write-invalidate}）．そのような書き込みが観測されれば，ロード
は古い値を読んでいたかもしれない．そのロードとそれ以降のすべての命令は
取り消され，分岐予測ミスの後とまったく同じように再実行される
（\cref{sec:l08-recovery}）．\margin{Intel のカウンタはこれをメモリ順序のマ
シンクリア（\texttt{MACHINE\_CLEARS}）と呼ぶ．1回につきパイプラインの再充
填，Skylake で約 17 サイクルを要する（第4章）．}ロードはまだコミットしていな
いので，これは常に
可能である．ロードがコミットしてしまえば，それはプログラム順に効力を生じて
おり，その位置をもはや監視する必要はない．

\begin{figure}[htbp]
  \centering
  % Sequential consistency with reordered loads.  Core 0 executes, in program
% order, I1: LW ready and I2: LW data; core 1 stores data (8) and then
% ready (1).  I2 executes early and reads the old data (3).  While I2 is
% uncommitted (shaded window) core 0 watches for writes to data by other
% cores; the invalidation caused by core 1's store makes it replay I2.
\begin{tikzpicture}[x=0.95cm,y=1cm, font=\footnotesize\sffamily, >={Stealth[length=1.6mm]},
  row/.style={ln-muted, line width=0.6pt},
  ev/.style={circle, fill=black, inner sep=1.3pt},
  lab/.style={font=\scriptsize, align=center, inner sep=1pt},
  win/.style={fill=ln-accent!20, draw=none}]
  % ---- row names
  \node[anchor=east] at (-0.15,1.4) {\iflangja{コア}{core}~0};
  \node[anchor=east] at (-0.15,0) {\iflangja{コア}{core}~1};
  % ---- monitoring windows of I2 (executed, not yet committed)
  \fill[win] (0.6,1.27) rectangle (2.6,1.53);
  \fill[win] (6.6,1.27) rectangle (9.8,1.53);
  \draw[row] (0,1.4) -- (10.5,1.4);
  \draw[row] (0,0) -- (10.5,0);
  % ---- core 0
  \node[ev] at (0.6,1.4) {};
  \node[lab, anchor=south] at (0.6,1.62) {\iflangja{I2: LW data\\3 を読む}{I2: LW data\\reads 3}};
  \node[ev, ln-caution] at (2.6,1.4) {};
  \node[lab, anchor=south, text=ln-caution] at (2.6,1.62) {\iflangja{I2 を\\再実行へ}{I2 must\\be replayed}};
  \node[ev] at (5.0,1.4) {};
  \node[lab, anchor=south] at (5.0,1.62) {\iflangja{I1: LW ready\\1 を読む}{I1: LW ready\\reads 1}};
  \node[ev] at (6.6,1.4) {};
  \node[lab, anchor=south] at (6.6,1.62) {\iflangja{I2: LW data\\8 を読む}{I2: LW data\\reads 8}};
  \node[ev] at (8.2,1.4) {};
  \node[lab, anchor=south] at (8.2,1.62) {\iflangja{I1\\コミット}{I1\\commits}};
  \node[ev] at (9.8,1.4) {};
  \node[lab, anchor=south] at (9.8,1.62) {\iflangja{I2\\コミット}{I2\\commits}};
  % ---- core 1
  \node[ev] at (2.6,0) {};
  \node[lab, anchor=north] at (2.6,-0.12) {\iflangja{SW data\\（8）}{SW data\\(8)}};
  \node[ev] at (4.1,0) {};
  \node[lab, anchor=north] at (4.1,-0.12) {\iflangja{SW ready\\（1）}{SW ready\\(1)}};
  % ---- invalidation seen by core 0
  \draw[->, ln-caution, dashed] (2.6,0.1) -- (2.6,1.28)
    node[pos=0.45, left, lab, text=ln-caution] {\iflangja{data の\\無効化}{invalidation\\of data}};
  % ---- time axis
  \draw[->, ln-muted] (0,-0.95) -- (10.5,-0.95)
    node[below left, inner sep=2pt, font=\scriptsize] {\iflangja{時間}{time}};
\end{tikzpicture}

  \caption{順序変更されたロードの再実行．プログラム順では，コア 0 は I1:
    \texttt{LW ready}，次に I2: \texttt{LW data} を実行する．I2 が実行済み
    でコミットしていない間（網掛け）にコア 1 がデータを書き，その無効化に
    よって I2 は再実行される．}
  \label{fig:l21-replay}
\end{figure}

\cref{fig:l21-replay} は，これを \cref{tab:l21-ready} のデータ準備完了フラグ
に適用したものである．コア 0 はデータのロード I2 をフラグのロード I1 より
前に実行し，古い値 3 を読む．その後コア 1 がデータに 8 をストアする．無効化
は I2 が未コミットのうちにコア 0 へ届くので，I2 は再実行され，今度は 8 を
読む．コア 1 がフラグを立てた後に実行された I1 は 1 を読み，両方のロードは，
プログラム順に実行されたときに返したはずの値のままコミットする．

再実行が必要なのは，あるロードの実行とコミットの間に，他のコアがそのロード
が読む位置へ書いたときだけである．これは大半のコードでは稀だが，データ準備
完了フラグの読み出し側のように，他のコアが書く位置をポーリングするループ
では頻繁に起こる．無効化はブロック全体を指すので，同じブロック内の別の位置
への書き込みも再実行を引き起こす．これは時間を要するが，誤った結果を生む
ことは決してない．\margin{ブロックは x86 でも多くの Arm コアでも 64 バイトな
ので，無関係な変数が同じブロックを共有する．偽共有（第19章）
はコヒーレンスミスだけでなく再実行も引き起こす．}

\begin{keypoint}
  逐次一貫性はメモリアクセスの順序変更を禁じるのではなく，プログラムが
  観測できる順序変更を禁じる．アウトオブオーダコアは，コミットより前に
  他のコアがその位置へ書いたロードをすべて再実行するならば，ロードを早く
  実行してよい．
\end{keypoint}

\section{緩和一貫性}
\label{sec:l21-relaxed}

逐次一貫性は，プログラムがその順序に依存するかどうかにかかわらず，コアが
行うアクセスのあらゆる対を制約する．\source{Lesson 21，動画 9--11；H\&P
§5.6．}あるコアがプログラム順に行う2つのアクセス，1つ目が位置 A への，
2つ目が位置 B へのものについては，WR A $\to$ WR B，WR A $\to$ RD B，
RD A $\to$ WR B，RD A $\to$ RD B の4種類の順序付けがある．いずれも，すべて
のコアから見て1つ目のアクセスが2つ目より前に効力を生じることを要求する
（\cref{tab:l21-orderings}）．逐次一貫性は4つすべてを常に守る．
\margin{x86 が緩和するのは WR $\to$ RD だけである．Arm，POWER，RISC-V は4つ
すべてを緩和する．Intel SDM 第3A巻「Memory Ordering」．}

\begin{table}[htbp]
  \centering\small
  \caption{あるコアの2つのアクセス，最初が A へ，次が B へのものの間の
    4つの順序付けと，それらに違反するコアが起こしうること．}
  \label{tab:l21-orderings}
  \begin{tabular}{@{}lp{4.3cm}p{4.1cm}@{}}
    \toprule
    順序付け & 違反 & 例 \\
    \midrule
    WR A $\to$ WR B & A の新しい値より前に B の新しい値が見える
      & データがストアされる前にフラグが立つ \\
    WR A $\to$ RD B & A への書き込みが見えるようになる前に B が読まれる
      & 2つのフラグ（\cref{ex:l21-flags}） \\
    RD A $\to$ WR B & A が読まれる前に B の新しい値が見える
      & 要求が読まれる前に応答が見える \\
    RD A $\to$ RD B & A より前に B が読まれる
      & データ準備完了フラグ（\cref{tab:l21-ready}） \\
    \bottomrule
  \end{tabular}
\end{table}

\begin{definition}[緩和一貫性]
  \term[かんわいっかんせい]{緩和一貫性}[relaxed consistency]を持つシステム
  は，通常のメモリアクセスについてこれらの順序付けのいくつかを常には守らず，
  正しさが順序に依存する箇所でプログラムが順序を強制する手段を提供する．
\end{definition}

読み出しで終わる2つの順序付け，WR $\to$ RD と RD $\to$ RD が，緩和の自然な
候補である．ロードが，同じコアの他の位置への先行するアクセスすべてより前に
効力を生じてよいならば，\cref{sec:l21-implementing} のコアは，コヒーレンス
トラフィックを監視する必要も，他のコアの書き込みのためにロードを再実行する
必要もない．どのロードも，\cref{fig:l21-reorder} や \cref{tab:l21-ready} の
ように，準備ができ次第実行するだけである．RD $\to$ RD だけを緩和しても，
ロードを監視しなければならない期間が短くなるにすぎない．ロードは依然として，
自分のコアの先行するストアすべての後に効力を生じなければならないからである．

\subsection{同期のためのアクセス}

2つの読み出しさえ順序変更されうるなら，プログラムは順序が重要な箇所で特定の
順序を得る手段を必要とする．そこで緩和一貫性を持つシステムは，2種類の操作を
区別する．通常のロードとストアは，モデルが許す範囲で順序変更されてよい．
順序が重要な箇所でプログラムが用いる特別な操作は，決して順序変更されない．
講義はそのような命令を \texttt{MSYNC}\index{MSYNC} と呼ぶ．プログラム順で
\texttt{MSYNC} に先行するすべてのメモリアクセスは，それに続くどのアクセス
よりも前に効力を生じる（\cref{fig:l21-msync}）．\margin{x86 にはそのような
命令 \texttt{MFENCE}\index{MFENCE} がある．その前のアクセスは，後に続く
どのアクセスよりも前に大域的に可視となる．}2つの \texttt{MSYNC} 命令の間
では，アクセスは依然として順序変更されうる．\caution{\texttt{MSYNC} は
\emph{完全な}フェンスである．実際の命令セットには部分的なフェン
ス---x86 の \texttt{SFENCE}，Arm の \texttt{DMB ISHST}---もあり，4つの対の一
部だけを順序付ける．}

\begin{figure}[htbp]
  \centering
  % Relaxed consistency with MSYNC.  The accesses of one core, in program
% order, are divided into groups by MSYNC instructions.  Accesses within a
% group may take effect in a different order; no access takes effect before
% an access that precedes it across an MSYNC.
\begin{tikzpicture}[x=1cm,y=1cm, font=\footnotesize\sffamily, >={Stealth[length=1.6mm]},
  acc/.style={draw, fill=ln-box, minimum width=1.05cm, minimum height=0.55cm, inner sep=1pt, font=\footnotesize\ttfamily},
  fence/.style={fill=ln-accent, minimum width=0.14cm, minimum height=1.1cm, inner sep=0pt},
  lab/.style={font=\scriptsize, align=center, inner sep=1pt}]
  % ---- accesses in program order
  \node[acc] (a1) at (0.0,0) {LW A};
  \node[acc] (a2) at (1.25,0) {SW B};
  \node[acc] (a3) at (2.5,0) {LW C};
  \node[fence] (f1) at (3.45,0) {};
  \node[acc] (a4) at (4.4,0) {LW D};
  \node[acc] (a5) at (5.65,0) {SW E};
  \node[acc] (a6) at (6.9,0) {LW F};
  \node[fence] (f2) at (7.85,0) {};
  \node[acc] (a7) at (8.8,0) {SW G};
  \node[acc] (a8) at (10.05,0) {LW H};
  \node[lab, anchor=south, text=ln-accent] at (f1.north) {MSYNC};
  \node[lab, anchor=south, text=ln-accent] at (f2.north) {MSYNC};
  % ---- reorderings allowed within a group
  \draw[<->, ln-tip] (a1.north) to[bend left=40] (a3.north);
  \draw[<->, ln-tip] (a4.north) to[bend left=40] (a6.north);
  \draw[<->, ln-tip] (a7.north) to[bend left=40] (a8.north);
  % ---- a reordering across MSYNC is not allowed
  \draw[->, ln-caution, dashed] (a5.south) to[bend left=35]
    node[pos=0.5, fill=white, inner sep=0.5pt, font=\small\bfseries] {$\times$} (a2.south);
  \node[lab, text=ln-caution, anchor=north] at (3.45,-1.15) {\iflangja{MSYNC をまたぐ\\順序変更は不可}{no reordering\\across MSYNC}};
  \node[lab, text=ln-tip, anchor=south] at (1.25,0.8) {\iflangja{グループ内は順序変更可}{reordering allowed}};
  \node[lab, text=ln-tip, anchor=south] at (5.65,0.8) {\iflangja{グループ内は順序変更可}{reordering allowed}};
  % ---- program order
  \draw[->, ln-muted] (-0.5,-1.95) -- (10.6,-1.95)
    node[below left, inner sep=2pt, font=\scriptsize] {\iflangja{プログラム順序}{program order}};
\end{tikzpicture}

  \caption{\texttt{MSYNC} は，コアのアクセスをプログラム順にグループへ
    分ける．1つのグループ内の異なる位置へのアクセスは，モデルが許す範囲で
    異なる順序で効力を生じてよい．前のグループのアクセスより前に効力を
    生じるアクセスはない．}
  \label{fig:l21-msync}
\end{figure}

両側に \texttt{MSYNC} を1つずつ置けば，データ準備完了フラグは緩和一貫性の
もとでも正しい．書き込み側と読み出し側は次を実行する．
\begin{lstlisting}[language={[x86masm]Assembler}, numbers=none]
; 書き込み側
      SW   R6, 0(R5)    ; data  <- 新しい値
      MSYNC             ; data が可視になった
      SW   R7, 0(R4)    ; ready <- 1
; 読み出し側
poll: LW   R1, 0(R4)    ; R1 <- ready
      BEQZ R1, poll     ; 未設定なら再度ポーリング
      MSYNC             ; ready のロードが完了した
      LW   R2, 0(R5)    ; R2 <- data
\end{lstlisting}
読み出し側では，\texttt{MSYNC} によって，フラグが立っているのを見たフラグ
のロードがデータのロードより前に効力を生じる．書き込み側では，データの
ストアがフラグへのストアより前に効力を生じるようにする．WR $\to$ WR が
緩和されていると，\texttt{MSYNC} がなければフラグへのストアの方が先に
見えてしまいかねないからである．読み出し側のフラグのロードが 1 を返した
ので，そのロードが効力を生じた時点でフラグは，そしてそれより前にデータは，
ストアされていた．データのロードはさらにその後に来る．読み出しで終わる
順序付けだけが緩和されるなら，書き込み側の2つのストアはすでにプログラム順
に効力を生じており，読み出し側の \texttt{MSYNC} だけで十分である．書き込み
側の \texttt{MSYNC} が必要になるのは，WR $\to$ WR も緩和されうる場合である．

フラグは同期の単純な形であり，2つの \texttt{MSYNC} の配置は一般的な規則に
従う．すなわち，先へ進む権利を獲得する操作の\emph{後}と，それを解放する操作
の\emph{前}に \texttt{MSYNC} を置く．\tip{2つのフェンスがクリティカルセク
ションの\emph{内側}を向いていると覚えるとよい．
どちらの端からも外へ出られない．}ロック（第20章）を使うスレッドは，
\texttt{lock} の直後に \texttt{MSYNC} を実行して，自分のクリティカル
セクションのアクセスが獲得の後に効力を生じ，それ以前のロック保持者の書き込み
をすべて見るようにし，\texttt{unlock} の直前にも実行して，次のスレッドが
ロックを獲得する時点で自分のクリティカルセクションの書き込みがすべて効力を
生じているようにする．このような命令は通常，あらゆるプログラムの中ではなく，
同期ライブラリのロックやバリアのルーチンの中に置かれる．
\margin{Arm はフェンスをアクセス自体に畳み込む．\texttt{LDAR}（獲得）
と \texttt{STLR}（解放）は一方向だけを順序付ける．RISC-V に
は \texttt{.aq}/\texttt{.rl} ビットがある．}

\section{データ競合と一貫性}
\label{sec:l21-races}

緩和一貫性がプログラムの結果を変えうるかどうかは，そのプログラムが共有
データにどうアクセスするかによる．\source{Lesson 21，動画 12；H\&P §5.6；
\cite{adve1990}．}

\begin{definition}[データ競合]
  異なるコアが行う同じメモリ位置への2つのアクセスは，少なくとも一方が
  書き込みであり---RD $\to$ WR，WR $\to$ RD または WR $\to$ WR---かつ同期
  によって順序付けられていない，すなわちどちらが先に効力を生じるかを定める
  同期操作が存在しないとき，\term[でーたきょうごう]{データ競合}[data race]%
  をなす．
\end{definition}

\begin{figure}[htbp]
  \centering
  % Data race versus accesses ordered by synchronization.  Left: a load and a
% store of A on different cores with nothing between them; either may take
% effect first.  Right: a barrier after the load on core 0 and before the
% store on core 1 forces the load to take effect first.
\begin{tikzpicture}[x=1cm,y=1cm, font=\footnotesize\sffamily, >={Stealth[length=1.6mm]},
  acc/.style={draw, fill=ln-box, minimum width=1.5cm, minimum height=0.5cm, inner sep=1pt, font=\footnotesize\ttfamily},
  bar/.style={draw, fill=ln-accent!20, minimum width=1.5cm, minimum height=0.5cm, inner sep=1pt, font=\scriptsize\ttfamily},
  hd/.style={font=\scriptsize, text=ln-muted},
  ttl/.style={font=\footnotesize\sffamily\bfseries}]
  % ================= left: data race
  \node[hd] at (0,0.6) {\iflangja{コア}{core}~0};
  \node[hd] at (3,0.6) {\iflangja{コア}{core}~1};
  \node[acc] (l0) at (0,0) {LW A};
  \node[acc] (s1) at (3,-1.1) {SW A};
  \draw[<->, ln-caution, line width=0.8pt] (l0.east) to[bend left=25] (s1.north);
  \node[font=\scriptsize, align=center, text=ln-caution] at (0.75,-1.1)
    {\iflangja{どちらが先かは\\タイミング次第}{either order,\\depending on timing}};
  \node[ttl, text=ln-caution] at (1.5,-2.5) {\iflangja{データ競合}{data race}};
  % ================= right: ordered by a barrier
  \begin{scope}[xshift=6.2cm]
    \node[hd] at (0,0.6) {\iflangja{コア}{core}~0};
    \node[hd] at (3,0.6) {\iflangja{コア}{core}~1};
    \node[acc] (l0b) at (0,0) {LW A};
    \node[bar] (b0) at (0,-0.8) {BARRIER};
    \node[bar] (b1) at (3,-0.8) {BARRIER};
    \node[acc] (s1b) at (3,-1.6) {SW A};
    \draw[ln-accent, densely dashed] (b0.east) -- (b1.west);
    \draw[->, ln-tip, line width=0.8pt] (l0b.south east) -- (s1b.north west);
    \node[ttl, text=ln-tip] at (1.5,-2.5) {\iflangja{データ競合なし}{no data race}};
  \end{scope}
  % ---- time runs downwards
  \draw[->, ln-muted] (-1.1,0.3) -- (-1.1,-2.0)
    node[below, font=\scriptsize, inner sep=2pt] {\iflangja{時間}{time}};
\end{tikzpicture}

  \caption{左: A のロードとストアを順序付けるものが何もない，データ競合．
    右: ロードの後とストアの前に置かれたバリアによって，ロードが先に効力を
    生じる．}
  \label{fig:l21-race}
\end{figure}

\cref{fig:l21-race} の左では，コア 0 の A のロードとコア 1 の A へのストア
の間には何もない．どちらが先に効力を生じるかはコアのタイミングによるので，
ロードはそれに応じて古い値か新しい値を返す．右では，コア 0 は自分のロードの
後にバリア（\cref{sec:l20-barrier}）に到達し，コア 1 は自分のストアの前に
同じバリアに到達する．すべてのスレッドが到達するまでどのスレッドもバリアを
離れないので，ロードは常にストアより前に効力を生じ，データ競合は存在しない．
同じロックで保護されたアクセスも順序付けられる．一方のクリティカルセクション
が \texttt{unlock} で終わってから，他方が \texttt{lock} で始まるからである．
\margin{競合検出器（ThreadSanitizer，Helgrind）が探すのはまさにこれ，
すなわちロックもバリアもフェンスも順序付けていない2つの衝突するアクセスであ
る．}

データ競合が重要なのは，次の性質による \cite{adve1990}．
\begin{itemize}
  \item \textbf{データ競合のないプログラム}は，その同期を認識するどの一貫性
    モデルのもとでも，逐次一貫性のもとでの結果と同じ結果を生む．同期を
    認識するモデルとは，同期操作自体が逐次一貫性を持ち，\texttt{MSYNC} が
    そうであるように，各同期操作が同じコアの先行するアクセスと後続の
    アクセスを隔てるものである．異なる位置へのアクセスや同じ位置の読み出し
    の順序を変えても，読まれる値はどれも変わらず，順序が変われば値が変わり
    うるアクセスの対はすべて同期によって順序付けられていて，そのような
    モデルはそれを順序どおりに保つからである．
  \item \textbf{データ競合のあるプログラム}には，そのような保証はない．
    緩和一貫性を持つシステムでは，どのインターリーブでも説明できない結果も
    含め，何でも起こりうる．
\end{itemize}

したがってプログラムは，まず逐次一貫性のもとで開発しデバッグできる．
そこでは，データ競合のあるプログラムでも，何らかのインターリーブで説明
できる結果を生むからである．同期が競合を取り除いた後は，速度のために
緩和一貫性のもとで実行すればよい．\sidenote{プログラミング言語はこれを契約に
した．C++11 と Java は，データ競合のないプログラムが逐次一貫性のもとと同じよ
うに振る舞うことを保証し，C++ ではデータ競合があるとプログラム全体が未定義に
なる．これらのモデルが認識する同期は \texttt{std::atomic}---既定
は \texttt{memory\_order\_seq\_cst} で，より安価な release/acquire も選べ
る---と，その上に築かれたロックである \cite{boehm2008}．}

この章の例はいずれもデータ競合を含む．\cref{ex:l21-flags} では，各コアが，
相手のコアが同期なしに書くフラグを読む．\cref{tab:l21-ready} のデータ準備
完了フラグは，メモリシステムが認識しない形の同期である．逐次一貫性のもとで
は，このフラグが読み出し側のデータのロードを書き込み側のストアより後に順序
付けるが，緩和一貫性を持つシステムはそれを知らない．\cref{sec:l21-relaxed}
の \texttt{MSYNC} はその順序付けをシステムから見えるようにするので，データ
のロードはもはやストアと競合しない．フラグをポーリングするロード自体は同期
されないままだが，フラグは 0 から 1 へしか変化しないので，それらは読み出し側
を遅らせうるだけである．\cref{tab:l20-race} の失われたインクリメントもデータ
競合である．逐次一貫性のもとでは，起こりうる結果は少なくとも2つのスレッドの
インターリーブが与えるものに限られる．緩和一貫性のもとでは，その限界すら
消える．\caution{逐次一貫性を持つシステムで動作するプログラムでも，データ
競合を含んでいれば緩和一貫性のもとでは失敗しうる．}

\begin{keypoint}
  同期をメモリシステムが認識する，データ競合のないプログラムは，緩和一貫性
  と逐次一貫性を見分けられない．したがって，一貫性を緩和しても，正しく同期
  されたプログラムは何も失わない．データ競合のあるプログラムは何でも観測
  しうる．
\end{keypoint}

\section{一貫性モデル}
\label{sec:l21-models}

逐次一貫性と \cref{sec:l21-relaxed} の緩和一貫性は，広い範囲にわたる一貫性
モデルの2つの点にすぎない．\source{Lesson 21，動画 13--14；
\cite{mosberger1993}．}モデルは，どの順序付けを緩和するか，そして同期を他の
アクセスと区別するかどうか，また区別するならどう区別するかで異なる．すべて
のアクセスを同じに扱うモデルもある．強い方から弱い方へ順に挙げる．
\begin{description}
  \item[強い一貫性] どの読み出しも，実時間で最も新しく書かれた値を返す
    （\cref{sec:l21-sc}）．
  \item[逐次一貫性] 各コアのプログラム順を保つ何らかのインターリーブが，
    すべての結果を説明する．
  \item[因果一貫性] \term[いんがいっかんせい]{因果一貫性}[causal
    consistency]のもとでは，因果的に関係する書き込み---同じコアの2つの
    書き込み，または，直接あるいは間接に先行する書き込みの結果を読んだ
    コアが行う書き込み---を，すべてのコアが同じ順序で見なければならない．
    並行な書き込みは，コアによって異なる順序で見えてよい．
  \item[プロセッサ一貫性]\index{ぷろせっさいっかんせい@プロセッサ一貫性}%
    あるコアの書き込みは，すべてのコアがそのコアの行った順序で見るが，
    異なるコアの書き込みを異なるコアが異なる順序で見ることがあり，また，
    あるコアの読み出しが，自分の他の位置への先行する書き込みより前に効力を
    生じることもある．
\end{description}
因果一貫性もプロセッサ一貫性も \cref{fig:l21-reorder} の実行を禁じない．
2つの書き込みは並行なので，各コアは相手のコアの書き込みを自分の読み出しの
後にはじめて見てよいからである．

他のモデルは同期を区別しており，それらが認識する同期を用いるデータ競合の
ないプログラムは，そのもとで逐次一貫性のもとと同じように振る舞う．
\begin{description}
  \item[弱い一貫性] \term[よわいいっかんせい]{弱い一貫性}[weak
    consistency]のもとでは，ロックのような同期変数へのアクセスは逐次一貫性
    を持ち，\texttt{MSYNC} のように振る舞う．すなわち，先行するアクセスは
    すべてそれより前に効力を生じ，後続のアクセスはどれもそうならない．
    その他のアクセスは，それらの間で自由に順序変更されてよい．
  \item[リリース一貫性]\index{りりーすいっかんせい@リリース一貫性}%
    同期は，\texttt{lock} と \texttt{unlock} のように，獲得と解放に分けられ
    る．獲得は後続のアクセスがそれより前に効力を生じることを妨げ，解放は
    先行するアクセスがすべて効力を生じるまで待つ．そのため通常のアクセスが
    クリティカルセクションの境界を越えて順序変更されることはない---
    \cref{sec:l21-relaxed} の \texttt{MSYNC} の配置である．
    \margin{リリース一貫性は，1990 年にスタンフォードの DASH マルチプロセッ
    サのために定義された．獲得と解放の分離をハードウェアの概念にした機械であ
    る．}
  \item[遅延リリース一貫性]\index{ちえんりりーすいっかんせい@遅延リリース一貫性}%
    リリース一貫性を遅延的に実現したもので，もとはソフトウェアで構築された
    分散共有メモリのためのものである．クリティカルセクションで行われた更新
    は，解放のたびにではなく，他のノードがロックを獲得してそのデータに
    アクセスするときにはじめて，そのノードへ転送される．
    \sidenote{遅延リリース一貫性 \cite{keleher1992} は，
    ワークステーションのネットワーク上のソフトウェア分散共有メモリのために設
    計された．そこでは更新の伝播はメッセージの送信を意味する．TreadMarks シ
    ステムは，ロックを獲得したノードがページに触れてはじめて，
    そのページを転送する．}
  \item[スコープ一貫性]\index{すこーぷいっかんせい@スコープ一貫性}%
    各ロックがスコープを定める．クリティカルセクションに入るスレッドは，
    同じロックの以前のクリティカルセクションで行われた更新を見るが，他の
    ロックのもとで行われた更新を必ずしも見ない \cite{iftode1996}．
\end{description}
この一覧は完全ではなく，新しいモデルが提案され続けている．それらに共通する
のは，この章の枠組みである．すなわち，システムはモデルが順序付けないアクセス
を順序変更してよく，特定の順序を必要とするプログラムは同期によってそれを
得る．

\needspace{18\baselineskip}
\begin{keypoint}[章のまとめ]
  コヒーレンスは1つの位置へのアクセスを順序付ける．メモリ一貫性が関わるのは，
  アウトオブオーダ実行と分岐予測が変えてしまう，異なる位置へのアクセスの
  順序である．逐次一貫性は，どの実行も，各コアのアクセスがプログラム順に
  行われ，異なるコアのアクセスがインターリーブされたかのように見えることを
  要求する．アクセスを一度に1つずつ行えばこれを提供できるが，アウトオブ
  オーダ実行を台無しにする．その代わりに，コアはロードを早く実行し，
  コミットより前に他のコアがその位置へ書いたロードを再実行すればよい．
  緩和一貫性はアクセス間の順序付けを常には保たず，獲得の後と解放の前に置く
  \texttt{MSYNC} のような同期操作をプログラムに与える．データ競合---同期に
  よって順序付けられていない，異なるコア上の衝突するアクセス---のない
  プログラムは，その同期を認識するどのモデルのもとでも，逐次一貫性のもとと
  同じように振る舞う．データ競合のあるプログラムは何でも観測しうる．弱い
  一貫性，リリース一貫性，遅延リリース一貫性，スコープ一貫性は同期の認識の
  仕方が異なり，因果一貫性とプロセッサ一貫性は，すべてのコアが一致しなければ
  ならない書き込みの順序が異なる．
\end{keypoint}

\end{document}
